\documentclass[preview]{standalone}

\usepackage{amsmath}
\usepackage{amssymb}
\usepackage{stellar}
\usepackage{definitions}

\begin{document}

\id{probability-exercises-batch2}
\genpage

\section{Exercises}

\begin{snippetexercise}{probability-exercises-2-ex-11}{}
    Let \(X\) and \(Y\) be two independent random variables with the same
    distribution, namely
    \[
        P(X=1) = P(X=-1) = \frac12,
    \]
    and define the random variable \(Z \triangleq XY\). Show that the three random
    variables are pairwise independent, but they are not independent.
\end{snippetexercise}

\begin{snippetsolution}{probability-exercises-2-ex-11-sol}{}
    \todo
\end{snippetsolution}

\begin{snippetexercise}{probability-exercises-2-ex-12}{}
    Let \(X\) and \(Y\) be two independent and identically distributed random
    variables with the following discrete density
    \[
        p(x) = \frac{1}{2^x}\chi_{\mathbb{N}\setminus\{0\}}(x).
    \]
    Show that \(p\) is a discrete density and compute:
    \begin{enumerate}
        \item \(P(\min(X,Y) \leq x)\);
        \item \(P(Y > X)\);
        \item \(P(X = Y)\);
        \item \(P(X \geq kY)\);
        \item \(P(\{X \text{ is a divisor of } Y\})\).
    \end{enumerate}
\end{snippetexercise}

\begin{snippetsolution}{probability-exercises-2-ex-12-sol}{}
    \todo
\end{snippetsolution}

\begin{snippetexercise}{probability-exercises-2-ex-13}{}
    Let \(X_1\), \(X_2\) and \(X_3\) be three independent random variables such that
    \(X_i \sim \operatorname{Geo}(1-p_i)\).
    \begin{enumerate}
        \item Show that
        \[
            P(X_1 < X_2 < X_3)
            =
            \frac{
                (1-p_1)(1-p_2)p_2p_3^2
            }{
                (1-p_2p_3)(1-p_1p_2p_3)
            }.
        \]
        \item Compute \(P(X_1 \leq X_2 \leq X_3)\).
    \end{enumerate}
\end{snippetexercise}

\begin{snippetsolution}{probability-exercises-2-ex-13-sol}{}
    \todo
\end{snippetsolution}

\begin{snippetexercise}{probability-exercises-2-ex-14}{}
    Let \(X\) be a discrete random variable with probability distribution
    \[
        p_k = c\frac{2^{-k}}{k}, \quad k \geq 1.
    \]
    \begin{enumerate}
        \item Show that
        \[
            \log x = \sum_{k\geq 1}
            (-1)^{k+1}\frac{(x-1)^k}{k}.
        \]
        \item Show that
        \[
            \log(x)
            =
            \log\left(
                \frac{
                    1 + \frac{x-1}{x+1}
                }{
                    1 - \frac{x-1}{x+1}
                }
            \right),
        \]
        and use this to prove that
        \[
            \log(x)
            =
            2\sum_{k\geq 1}
            \frac{1}{2k-1}
            \left(
                \frac{x-1}{x+1}
            \right)^{2k-1}.
        \]
        \item Determine for which constant \(c\) the sequence
        \((p_k)_{k\geq 1}\) is a probability distribution.
        \item Compute the mode of the distribution \((p_k)_{k\geq 1}\), that is,
        the value at which the distribution attains its maximum.
        \item Compute the probabilities of the events
        \[
            A = \{X > 1\}
            \quad\text{and}\quad
            B = \{X \text{ takes even values}\}.
        \]
    \end{enumerate}
    \emph{Hint:} use the fact that
    \[
        \log(1+y)
        =
        \sum_{k\geq 1}
        (-1)^{k+1}\frac{y^k}{k}.
    \]
\end{snippetexercise}

\begin{snippetsolution}{probability-exercises-2-ex-14-sol}{}
    \todo
\end{snippetsolution}

\begin{snippetexercise}{probability-exercises-2-ex-15}{}
    A coin is tossed \(30\) times, and \(S_n\) denotes the number of successes in
    the first \(n\) trials.
    \begin{enumerate}
        \item Compute the probability that in the \(30\) tosses there are \(n\)
        successes, knowing that in the first \(10\) tosses only successes were
        obtained.
        \item Compute the probability that in the \(30\) tosses there are \(n\)
        successes, knowing that at least \(10\) successes were obtained.
    \end{enumerate}
\end{snippetexercise}

\begin{snippetsolution}{probability-exercises-2-ex-15-sol}{}
    \todo
\end{snippetsolution}

\begin{snippetexercise}{probability-exercises-2-ex-16}{}
    Two random variables \(X_1\) and \(X_2\) are independent and uniformly
    distributed, respectively, on \(\{0,1,2\}\) and on \(\{0,1,2,3\}\). Compute the
    distribution of the transformation
    \[
        Z = |X_1 - X_2|.
    \]
\end{snippetexercise}

\begin{snippetsolution}{probability-exercises-2-ex-16-sol}{}
    \todo
\end{snippetsolution}

\begin{snippetexercise}{probability-exercises-2-ex-17}{}
    Let \((X,Y)\) be a random vector taking values in
    \[
        E \triangleq \{1,2,3,4\}^2
    \]
    with distribution
    \[
        P_{(X,Y)}((i,j))
        =
        \frac{C}{i(5-j)},
        \quad C \in \realnumbers,
        \quad (i,j)\in E.
    \]
    \begin{enumerate}
        \item Determine for which value of the constant \(C\) the function
        \(P_{(X,Y)}\) is a probability distribution.
        \item Check whether the two random variables are independent.
        \item Compute
        \[
            P(Y \in \{2,3\}\mid X=2).
        \]
        \item Compute the covariance of the vector.
    \end{enumerate}
\end{snippetexercise}

\begin{snippetsolution}{probability-exercises-2-ex-17-sol}{}
    \todo
\end{snippetsolution}

\begin{snippetexercise}{probability-exercises-2-ex-18}{}
    \(X\) and \(Y\) are two independent random variables with uniform distribution
    on \(\{0,1,2,3,4\}\). Compute the probability that the equation
    \[
        t^2 + X \cdot t + Y = 0
    \]
    in the variable \(t\) has two distinct real solutions. What is the probability
    that it has a unique solution?
\end{snippetexercise}

\begin{snippetsolution}{probability-exercises-2-ex-18-sol}{}
    \todo
\end{snippetsolution}

\begin{snippetexercise}{probability-exercises-2-ex-19}{}
    Let \((X,Y)\) be a random vector with values in the set
    \[
        S = \{(k,j) \suchthat k,j\in\mathbb{N}\cup\{0\},\, j\geq k\}
    \]
    with the following probability distribution
    \[
        p(k,j)
        \triangleq
        \begin{cases}
            0 & j < k, \\
            \frac{
                e^{-(\lambda+\mu)}\lambda^k\mu^{j-k}
            }{
                k!(j-k)!
            } & j \geq k.
        \end{cases}
    \]
    \begin{enumerate}
        \item Compute the marginal distributions of the random vector.
        \item Compute, if it exists finite, the covariance matrix of the random
        vector.
        \item Are \(X\) and \(Y\) independent?
    \end{enumerate}
\end{snippetexercise}

\begin{snippetsolution}{probability-exercises-2-ex-19-sol}{}
    \todo
\end{snippetsolution}

\begin{snippetexercise}{probability-exercises-2-ex-20}{}
    There are two urns \(U_1\) and \(U_2\), containing balls marked \(0\) and \(1\).
    The composition of the urns is
    \[
        U_1 = \{n_{1,0}, n_{1,1}\},
        \quad
        U_2 = \{n_{2,0}, n_{2,1}\},
    \]
    where \(n_{i,j}\) is the number of balls marked \(j\) in urn \(i\).
    A ball is drawn from the first urn and transferred to the second urn. Then a
    ball is drawn from the second urn.
    \begin{enumerate}
        \item Construct the probability space associated with the draw from the
        second urn.
        \item Compute the probability that the ball drawn from the first urn was
        marked \(1\), knowing that a ball marked \(1\) was drawn from the second urn.
        \item Let \(X_1\) and \(X_2\) be the results of the draws from the first and
        the second urn. Check whether the two random variables are independent.
        \item Compute the marginal probability distributions of \(X_1\) and \(X_2\).
        \item Compute the distribution of \(S \triangleq X_1 + X_2\).
        \item Compute the probability that a \(1\) was obtained in the first draw,
        knowing that \(S=2\).
    \end{enumerate}
\end{snippetexercise}

\begin{snippetsolution}{probability-exercises-2-ex-20-sol}{}
    \todo
\end{snippetsolution}

\begin{snippetexercise}{probability-exercises-2-ex-21}{}
    Let \(X\) be a discrete random variable with distribution symmetric with
    respect to \(0\), and define the random variable
    \[
        Y = X^{2k}
    \]
    for a fixed \(k\in\mathbb{N}\). Assuming that
    \[
        \mathbb E\left[|X|^{2k+1}\right] < \infty,
    \]
    \begin{enumerate}
        \item compute \(\mathbb E[X]\) and \(\mathbb E[Y]\);
        \item compute \(\mathbb E[XY]\);
        \item are the two random variables independent?
    \end{enumerate}
\end{snippetexercise}

\begin{snippetsolution}{probability-exercises-2-ex-21-sol}{}
    \todo
\end{snippetsolution}

\end{document}